%% Q1
A doctor serves and treats a patient; a \textbf{teacher} serves and teaches a \textbf{student}. The relationship is \emph{professional : the person served}.

%% Q2
Apple, orange and banana are \emph{kinds} of fruit; \textbf{fruit} is the general class, so it is different.

%% Q3
ACFJ: A(1), C(3), F(6), J(10), with gaps $+2, +3, +4$. OUBJ continues the same letters: O(15, $+5$), U(21, $+6$), B(28$\to$2, $+7$), J(10, $+8$). So the second group continues the first.
SUXB: S(19), U(21), X(24), B(28$\to$2). Continue: G(7, $+5$), M(13, $+6$), T(20, $+7$), B(28$\to$2, $+8$), giving \textbf{GMTB}.

%% Q4
BEAT $\rightarrow$ GIDV: B$+5$, E$+4$, A$+3$, T$+2$.
SOUP $\rightarrow$ S$+5=$X, O$+4=$S, U$+3=$X, P$+2=$R, giving \textbf{XSXR}.

%% Q5
Remaining calm and \textbf{silently looking at the students} usually restores order without confrontation. Shouting or scolding adds to the noise.

%% Q6
A teacher's job comprises \textbf{propagation, collection (preservation) and creation of knowledge}, i.e.\ teaching, scholarship and research together.

%% Q7
A good researcher is \textbf{open-minded and radical}: willing to question established ideas and accept new evidence.

%% Q8
An effective teacher \textbf{motivates students to learn}.

%% Q9
\textbf{Symbiosis} (mutualism) is an association in which \textbf{both partners benefit}, e.g.\ lichens (alga $+$ fungus). (The options loosely call them host and parasite.)

%% Q10
HIV, the virus that causes AIDS, infects \textbf{helper T-lymphocytes (CD4 cells)}, which are essential for triggering antibody production and immune responses. This destroys the body's immunity.

%% Q11
\textbf{Pasteurisation} (Louis Pasteur) heats milk (e.g.\ $72^\circ$C for 15 s, then rapid cooling) to \textbf{destroy harmful micro-organisms}.

%% Q12
Genetic engineering means deliberately modifying genes \textbf{using recombinant DNA (rDNA) technology}: cutting and joining DNA from different sources.

%% Q13 ctx
\textbf{Read the following passage and answer question nos.~13 to 16.}\par\smallskip
From the very beginning man has attempted what has seemed impossible. Man is different from the rest of the creation in this respect. He has an eternal thirst for adventure. This has led to countless new discoveries and inventions. Human curiosity is limitless. It is responsible for space flights and moon landings.

The desire to know what lies beyond the visible world takes many forms. The everest hero Tenzing and the hero of the seven Seas, Mihir Sen, were inspired by the same restless spirit. Astronaut Armstrong and his colleagues, who were the first humans to set foot on the soil of the moon, have proved beyond doubt that man shall not rest until he has conquered the entire universe.

Now the question arises: Is it enough to know and master nature? To answer it one should first answer a more basic question: which is more important --- knowing and understanding the world around you or knowing and understanding yourself? Everyone will agree that in the absence of self-knowledge even the most advanced knowledge of the universe will be not only useless but also dangerous.

%% Q13
``Human curiosity is limitless. \textbf{It} is responsible for space flights and moon landings.'' \emph{It} refers to \textbf{man's curiosity}.

%% Q14
``Man shall not rest'' means man shall not \textbf{stop}.

%% Q15
Per the key, ``eternal thirst'' is a thirst \textbf{that cannot be satisfied} (a never-ending, insatiable thirst). (Literally, \emph{eternal} means without beginning or end.)

%% Q16
``Man is different from the rest of the creation'' refers to \textbf{all mankind} (the human species).

%% Q17 ctx
\textbf{Read the following passage and answer question nos.~17 to 20.}\par\smallskip
The word medieval means `the middle age'. It is used in history, therefore, to refer to that period which lies between the ancient and the modern periods and is quite literally the middle period. How do we know when the ancient period ends and the medieval period begins? We have taken the eighth century A.D.\ as the beginning and the eighteenth century A.D.\ as the end of the medieval period. Why? Because as you will see when you read this book, there were many changes taking place in Indian Society in about the eighth century A.D.\ and during this time these changes influenced many aspects of Indian Society. They influenced the political and economic aspects, social laws, religion, language, art in short, almost everything. So we recognise that a new stage in Indian history had come about. We may say that this change took place around the eighth century.

We have to be just a little vague about the date because just as a person does not change suddenly in one year, similarly societies also take time to change. Nor does everybody become influenced by new ideas at the same time. Some of these changes in Indian history began earlier than the eighth century and in some parts of India their influence was felt a little later. But taking a general view of things, we can say that the new phase took root in the eighth century. In the same way with the break-up of the Mughal Empire and the coming of the British, the eighteenth century also saw many changes. So we refer to the century as the closing of the medieval period.

%% Q17
``Vague about the date'' means \textbf{not precise}: the date cannot be fixed exactly, because change is gradual.

%% Q18
The author says that around the eighth century ``a new stage in Indian history had come about'', which ended with the changes of the eighteenth century. So \textbf{a new stage began in the eighth century and lasted up to the eighteenth}.

%% Q19
``Quite literally the middle period'' means the word, taken in its exact sense, means ``middle period''. \emph{Literally} here means \textbf{accurately} (exactly).

%% Q20
``\ldots as you will see when you \textbf{read this book}'' shows the extract is from \textbf{a book}.

%% Q21
Personality develops from the interaction of \textbf{heredity, environment and society}: all three.

%% Q22
\textbf{Amnesia} is \textbf{loss of memory}.

%% Q23
Long-term memory stores information over long periods, e.g.\ \textbf{remembering an event of the long past}.

%% Q24
Tax $=15\%$ of salary $=$ Rs.~30,000, so salary $=30{,}000\div0.15=$ \textbf{Rs.~2,00,000}.

%% Q25
$\tfrac25a=\tfrac23b\Rightarrow\dfrac ab=\dfrac{2/3}{2/5}=\dfrac53$. The ratio is \textbf{5 : 3}.

%% Q26
Area to paint $=20\times14-7\times4=280-28=252$ sq ft. Cost $=252\times10=$ \textbf{Rs.~2,520}.

%% Q27
Sum $=4\times12+2\times15=48+30=78$; average $=78\div6=\mathbf{13}$.

%% Q28
31 December 2001 is Monday, so 1 January 2002 is Tuesday. 2002 has 365 days ($=52$ weeks $+1$ day), so 1 January 2003 is one day later: \textbf{Wednesday}.

%% Q29
Priya $=300$, Kanu $=200$, so Geeta $=600-500=$ \textbf{Rs.~100}.

%% Q30
Net cost $=400\times0.85=$ \textbf{Rs.~340}.

%% Q31
\textbf{Satyagraha} (``truth-force'') is Gandhi's method of non-violent resistance: persuading the opponent through truth, suffering, dialogue and negotiation rather than force. Among the options, \textbf{discussion and negotiation} is closest (per the key).

%% Q32
The chain is attraction $\rightarrow$ infatuation $\rightarrow$ anger $\rightarrow$ frustration. By the same chain, \textbf{infatuation causes frustration} (via anger). The other options reverse the direction of causation.

%% Q33
The argument is fallacious (Einstein does not become a reformer merely by living in the same century). The only correct statement among the options is the given premise, \textbf{``Gandhi was a social reformer''}.

%% Q34
P and S are brothers and M is P's son, so \textbf{S is M's uncle}. (T is the mother of S, not the grandmother; P is a brother, not a sister; M and S are nephew and uncle.)

%% Q35
Visible light spans about $4\times10^{-7}$ m to $7\times10^{-7}$ m (400--700 nm). None of the ranges given matches, so \textbf{none of the above}.

%% Q36 note
Mercury is a \textbf{metal} (the only metal that is liquid at room temperature), which is option (B). The official key gives (A); answer (B) in your own understanding.

%% Q36
Mercury (Hg) is an element, a \textbf{metal} that is liquid at room temperature. It is not a compound. The official key's letter (A) appears to be a misprint.

%% Q37 note
The kilowatt-hour is a unit of \textbf{energy} (option B). The official key gives (A).

%% Q37
$1\ \text{kWh}=1000\ \text{W}\times3600\ \text{s}=3.6\times10^6$ J. It is a unit of \textbf{energy} (the ``unit'' of electricity billing), not of power. The official key's (A) appears to be a misprint.

%% Q38
\textbf{Tidal energy} is a non-conventional (renewable) source. Wood, coal and mineral oil are conventional.

%% Q39
Naphtha is a light petroleum fraction used as a feedstock and solvent in the chemical and petrochemical industry. The official key gives \textbf{(B) chemical in a plant}.

%% Q40 note
A body lifted to a height gains \textbf{potential energy} ($mgh$), option (B). The official key gives (A).

%% Q40
Work done in lifting a body against gravity is stored as \textbf{gravitational potential energy} ($E_p=mgh$). The key's (A) Kinetic appears to be a misprint.

%% Q41
The \textbf{National Missile Defense} programme was pushed forward by President \textbf{George W.~Bush (Jr.)}, who withdrew from the ABM Treaty in 2001--02.

%% Q42
Economic liberalisation (1991) was initiated by the \textbf{P.~V.~Narasimha Rao} government, with Dr.~Manmohan Singh as Finance Minister.

%% Q43
India's first nuclear test, \emph{Smiling Buddha}, was at Pokhran on 18 May \textbf{1974}.

%% Q44
Bank nationalisation (1969) and the abolition of privy purses (26th Amendment, 1971) were initiated by \textbf{Indira Gandhi}.

%% Q45
\textbf{Sardar Vallabhbhai Patel}, as Home Minister, integrated over 560 princely states into the Indian Union.

%% Q46 note
Lord Mountbatten was the first Governor-General of independent India (1947--48); C.~Rajagopalachari was the first and only \emph{Indian} Governor-General (1948--50). The official key gives (C).

%% Q46
The key's intended answer is \textbf{C.~Rajagopalachari}, the first Indian to hold the office of Governor-General.

%% Q47
Let group C $=x$. Then B $=2x$ and A $=x/2$, so $x+2x+\tfrac x2=56\Rightarrow3.5x=56\Rightarrow x=\mathbf{16}$.

%% Q48
The setting sun is in the west. After turning left, the west is on your left, so you now face \textbf{north}. Before the left turn you were facing \textbf{east}.

%% Q49
The first duty is to save life: \textbf{render first-aid} to the injured child (then seek medical help and inform the police).

%% Q50
Mental ability is expressed as the \textbf{Intelligence Quotient (IQ)} $=\dfrac{\text{Mental age}}{\text{Chronological age}}\times100$.
